\documentclass[12pt,a4paper]{article}

\usepackage{iftex}
\ifPDFTeX
  \usepackage[T2A]{fontenc}
  \usepackage[utf8]{inputenc}
  \usepackage[russian]{babel}
\else
  \usepackage{fontspec}
  \usepackage[russian]{babel}
  \setmainfont{Arial Unicode MS}
  \setsansfont{Arial Unicode MS}
  \setmonofont{Courier New}
\fi

\usepackage{amsmath}
\usepackage{amssymb}
\usepackage{booktabs}
\usepackage{array}
\usepackage{enumitem}
\usepackage{geometry}
\usepackage{xcolor}
\usepackage{hyperref}

\geometry{left=18mm,right=18mm,top=18mm,bottom=18mm}
\hypersetup{colorlinks=true,linkcolor=teal,urlcolor=teal}
\setlist[itemize]{leftmargin=*,itemsep=2pt,topsep=4pt}
\renewcommand{\arraystretch}{1.18}

\newcommand{\code}[1]{\texttt{#1}}
\newcommand{\sectionline}{\vspace{0.4em}\noindent\textcolor{teal}{\rule{\linewidth}{0.6pt}}\vspace{0.4em}}
\newcommand{\simple}[1]{\par\smallskip\noindent\textbf{Простыми словами:} #1\par}
\newcommand{\howto}[1]{\par\smallskip\noindent\textbf{Как читать формулу:} #1\par}
\newcommand{\effect}[1]{\par\smallskip\noindent\textbf{Что влияет на результат:} #1\par}
\newcommand{\exampletext}[1]{\par\smallskip\noindent\textbf{Мини-пример:} #1\par}

\title{Формулы расчетной модели аэролодки}
\author{Цифровой штурман аэролодки}
\date{}

\begin{document}
\maketitle

\noindent
Документ описывает, как веб-приложение считает скорость, режим глиссирования,
мощность, расход топлива, риск и итоговую стоимость ребра графа.

\section{Общая идея модели}

Маршрут состоит из коротких участков графа. Для каждого участка известны длина,
тип поверхности, риск и источник ребра: широкий \code{navmesh}, узкая линия
\code{waterway} или привязка пользовательской точки. Дальше C++-движок
подбирает рекомендованную скорость, определяет режим движения и считает
стоимость участка.

\begin{itemize}
  \item Скорость зависит от выбранного режима маршрута.
  \item Глиссирование определяется через число Froude и минимальный порог скорости.
  \item Топливо считается через сопротивление, требуемую мощность и удельный расход двигателя.
  \item Риск увеличивается на быстрых и узких участках.
  \item Алгоритм A* выбирает путь с минимальной суммарной стоимостью.
\end{itemize}

\section{Основные переменные}

\begin{center}
\begin{tabular}{>{\ttfamily}p{0.18\linewidth}p{0.58\linewidth}p{0.16\linewidth}}
\toprule
\textnormal{Обозначение} & Смысл & Единицы \\
\midrule
d\_seg & длина участка графа & км \\
V\_rec & рекомендованная скорость на участке & км/ч \\
v & та же скорость в метрах в секунду & м/с \\
m\_dry & сухая масса аэролодки & кг \\
m\_payload & полезная загрузка & кг \\
m & полная масса & кг \\
L & расчетная длина корпуса & м \\
Fn & число Froude & безразмерно \\
R & оценка сопротивления движению & Н \\
P & требуемая мощность & кВт \\
BSFC & удельный расход топлива двигателя & г/кВт·ч \\
rho\_fuel & плотность топлива & кг/л \\
k\_surf & коэффициент поверхности & безразмерно \\
k\_mode & коэффициент выбранного режима & безразмерно \\
\bottomrule
\end{tabular}
\end{center}

\section{Как читать формулы, если математика дается тяжело}

В этом документе формула всегда отвечает на простой вопрос: какое число нужно
получить для одного участка маршрута. Не нужно заучивать все обозначения сразу.
Достаточно читать формулу слева направо:

\begin{itemize}
  \item слева от знака \(=\) написано, что мы хотим найти;
  \item справа от знака \(=\) написано, из каких уже известных чисел это считается;
  \item знак \(\cdot\) означает умножение;
  \item дробь означает деление;
  \item \(\max(a,b)\) означает ``возьми большее из двух чисел'';
  \item \(\sum\) означает ``сложи значения по всем участкам маршрута'';
  \item если написано ``иначе'', это значит: если условия выше не сработали.
\end{itemize}

Главная логика модели такая: сначала считаем скорость и режим движения, потом
сопротивление, затем мощность и топливо, а в конце риск и стоимость ребра для
алгоритма выбора маршрута.

\section{Формулы по шагам}

\subsection{Шаг 1. Полная масса}

\begin{align}
m &= \max\left(200,\; m_{\mathrm{dry}} + m_{\mathrm{payload}}\right), \\
m_{\mathrm{eff}} &= m \cdot k_{\mathrm{load}}.
\end{align}

Масса нужна для сопротивления поверхности. Чем тяжелее лодка, тем больше сила,
которую нужно преодолеть при движении по воде, льду, шуге или камышу.

\simple{Сначала складываем массу самой лодки и все, что положили внутрь:
людей, груз, топливо, оборудование. Затем умножаем на коэффициент конфигурации
\(k_{\mathrm{load}}\), потому что разные режимы лодки могут по-разному
нагружать корпус и винт.}
\howto{\(m_{\mathrm{dry}}\) -- вес пустой лодки, \(m_{\mathrm{payload}}\) --
загрузка. Если сумма вдруг получилась слишком маленькой, модель все равно берет
минимум 200 кг, чтобы не считать физически странные значения.}
\effect{Чем больше загрузка, тем больше \(m\) и \(m_{\mathrm{eff}}\).
Дальше из-за этого растет сопротивление, мощность и расход топлива.}
\exampletext{Лодка весит 1450 кг, загрузка 650 кг. Тогда
\(m=1450+650=2100\) кг. Если \(k_{\mathrm{load}}=1\), эффективная масса тоже
2100 кг.}

\subsection{Шаг 2. Перевод скорости и время участка}

\begin{align}
v &= \frac{V_{\mathrm{rec}}}{3.6}, \\
t_{\mathrm{seg}} &= \frac{d_{\mathrm{seg}}}{V_{\mathrm{rec}}}.
\end{align}

Время участка считается в часах. Если участок 1 км, а рекомендованная скорость
40 км/ч, время будет \(1/40\) часа, то есть 1.5 минуты.

\simple{Скорость в интерфейсе удобнее показывать в км/ч, но физические формулы
обычно работают с м/с. Поэтому скорость делится на 3.6. Потом считаем время:
длина участка делится на скорость.}
\howto{\(v\) -- скорость в м/с. \(t_{\mathrm{seg}}\) -- сколько часов лодка
идет по конкретному участку.}
\effect{Если скорость выше, время меньше. Если участок длиннее, время больше.}
\exampletext{Если \(V_{\mathrm{rec}}=36\) км/ч, то \(v=36/3.6=10\) м/с.
Если участок 2 км, то \(t_{\mathrm{seg}}=2/36\) часа, примерно 3.3 минуты.}

\subsection{Шаг 3. Число Froude}

\begin{equation}
Fn = \frac{v}{\sqrt{gL}}.
\end{equation}

Число Froude показывает отношение скорости лодки к характерной волновой
скорости корпуса.

\simple{Это число помогает понять, идет лодка медленно ``как корпус в воде''
или уже близка к глиссированию. Чем быстрее лодка и чем короче корпус, тем
больше \(Fn\).}
\howto{В числителе стоит скорость \(v\). В знаменателе стоит
\(\sqrt{gL}\): это условная скорость, связанная с длиной корпуса. Мы сравниваем
реальную скорость с этой условной скоростью.}
\effect{Если \(Fn\) маленькое, лодка скорее идет в водоизмещающем режиме. Если
\(Fn\) достаточно большое, лодка может выйти на глиссирование.}
\exampletext{Если \(v=12\) м/с, \(g=9.81\), \(L=6.9\), то
\(\sqrt{gL}\approx8.2\), а \(Fn\approx12/8.2\approx1.46\).}

\subsection{Шаг 4. Порог устойчивого глиссирования}

\begin{align}
m_{\mathrm{ref}} &= \max\left(m_{\mathrm{dry}} + 500,\; 1\right), \\
V_{\mathrm{planing}} &=
\max\left(
V_{\min},\;
3.6 \cdot Fn_{\mathrm{full}} \sqrt{gL}
\sqrt{\frac{m}{m_{\mathrm{ref}}}}
\right).
\end{align}

Порог глиссирования растет с массой. Тяжелая лодка требует большей скорости,
чтобы выйти на устойчивое скольжение.

\simple{Это минимальная скорость, после которой лодка может уверенно скользить,
а не толкать перед собой воду или рыхлую поверхность. Если лодка тяжелее, ей
нужна большая скорость для выхода на глиссирование.}
\howto{Сначала задаем базовую массу \(m_{\mathrm{ref}}\). Потом считаем
скорость по Froude и умножаем ее на поправку \(\sqrt{m/m_{\mathrm{ref}}}\).
Функция \(\max\) не дает порогу стать меньше минимальной скорости
\(V_{\min}\).}
\effect{Больше масса \(m\) -- выше порог. Больше длина корпуса \(L\) -- тоже
выше расчетная скорость. Больше \(Fn_{\mathrm{full}}\) -- строже условие
глиссирования.}
\exampletext{Если лодку сильно загрузить, то \(m/m_{\mathrm{ref}}\) станет
больше единицы. Корень из этого числа тоже станет больше единицы, значит
\(V_{\mathrm{planing}}\) вырастет.}

\subsection{Шаг 5. Выбор режима движения}

\begin{equation}
\mathrm{motion} =
\begin{cases}
\mathrm{planing}, &
\begin{aligned}
&\mathrm{surface.planing}=\mathrm{true},\ \neg\mathrm{narrow}, \\
&V_{\mathrm{rec}} \ge V_{\mathrm{planing}},\ Fn \ge Fn_{\mathrm{full}}, \\
&r_{\mathrm{surface}} \le 3,
\end{aligned}
\\[1.2em]
\mathrm{transition}, &
\begin{aligned}
&\mathrm{surface.planing}=\mathrm{true}, \\
&V_{\mathrm{rec}} \ge 0.72V_{\mathrm{planing}},\ Fn \ge Fn_{\mathrm{on}}, \\
&r_{\mathrm{surface}} \le 3,
\end{aligned}
\\[1.2em]
\mathrm{displacement}, & \text{иначе}.
\end{cases}
\end{equation}

На каждом участке движок помечает режим движения: глиссирование, переходный
режим или водоизмещающий режим.

\simple{Здесь модель выбирает один из трех ярлыков для участка. ``Planing'' --
лодка уверенно глиссирует. ``Transition'' -- лодка почти вышла на глиссирование,
но режим еще не самый эффективный. ``Displacement'' -- лодка идет тяжелее и
медленнее, без устойчивого скольжения.}
\howto{Формула читается как список условий. Если выполнены все условия первой
строки, выбираем глиссирование. Если первая строка не подошла, но подошла
вторая, выбираем переходный режим. Если не подошло ничего, выбираем
водоизмещающий режим.}
\effect{Узкий участок, большой риск поверхности или недостаточная скорость
мешают глиссированию. Поэтому безопасный режим часто переводит лодку в более
осторожное движение.}
\exampletext{Если поверхность разрешает глиссирование, участок широкий,
\(V_{\mathrm{rec}}\) выше порога и риск не больше 3, получаем
\(\mathrm{planing}\). Если скорость немного ниже уверенного порога, но уже
достаточная для начала выхода, получаем \(\mathrm{transition}\).}

\subsection{Шаг 6. Коэффициент сопротивления корпуса}

\begin{equation}
C_d =
\begin{cases}
C_{d,\mathrm{planing}}, & \mathrm{motion}=\mathrm{planing}, \\
\dfrac{C_{d,\mathrm{disp}} + C_{d,\mathrm{planing}}}{2},
& \mathrm{motion}=\mathrm{transition}, \\
C_{d,\mathrm{disp}}, & \mathrm{motion}=\mathrm{displacement}.
\end{cases}
\end{equation}

В глиссировании сопротивление корпуса ниже, в водоизмещающем режиме выше,
а переходный режим находится между ними.

\simple{\(C_d\) -- это ``насколько трудно корпусу идти через среду''. В
глиссировании корпус скользит легче, поэтому коэффициент меньше. В
водоизмещающем режиме корпусу тяжелее, поэтому коэффициент больше.}
\howto{Формула снова выбирает одно из трех значений. Для переходного режима
берется среднее между плохим и хорошим вариантом.}
\effect{Чем больше \(C_d\), тем больше сопротивление \(R\), а значит выше
мощность и расход топлива.}
\exampletext{Если \(C_{d,\mathrm{disp}}=0.42\), а
\(C_{d,\mathrm{planing}}=0.18\), то в transition получаем
\((0.42+0.18)/2=0.30\).}

\subsection{Шаг 7. Горб сопротивления}

\begin{equation}
k_{\mathrm{hump}} =
\begin{cases}
0.72, & \mathrm{motion}=\mathrm{planing}, \\
1.18, & \mathrm{motion}=\mathrm{transition}, \\
1 + 0.42\left(\dfrac{Fn}{\max(Fn_{\mathrm{on}},0.1)}\right)^4,
& \mathrm{motion}=\mathrm{displacement}.
\end{cases}
\end{equation}

Перед выходом на глиссирование лодка может попадать в неэффективную зону:
скорость уже высокая, но корпус еще не скользит устойчиво.

\simple{Горб сопротивления -- это момент, когда лодка уже разогналась, но еще
не вышла на нормальное скольжение. В этот момент топлива может уходить много,
а пользы от скорости мало.}
\howto{Если лодка уже глиссирует, множитель \(k_{\mathrm{hump}}\) меньше 1:
сопротивление снижается. В переходном режиме множитель больше 1:
сопротивление растет. В водоизмещающем режиме он зависит от \(Fn\).}
\effect{Большой \(k_{\mathrm{hump}}\) увеличивает сопротивление, а значит
увеличивает мощность и расход.}
\exampletext{Если \(k_{\mathrm{hump}}=1.18\), сопротивление динамической части
увеличится примерно на 18\%. Если \(k_{\mathrm{hump}}=0.72\), наоборот,
динамическая часть снизится примерно на 28\%.}

\subsection{Шаг 8. Сопротивление участка}

\begin{align}
R_{\mathrm{dyn}} &=
0.5 \rho_{\mathrm{air}} A_{\mathrm{res}} C_d v^2, \\
R_{\mathrm{air}} &=
0.5 \rho_{\mathrm{air}} A_{\mathrm{air}} \cdot 0.9 \cdot v^2, \\
R_{\mathrm{surf}} &=
m_{\mathrm{eff}}g\mu_{\mathrm{surface}}\max(0.55,k_{\mathrm{surf}}), \\
R &=
\left(R_{\mathrm{dyn}}k_{\mathrm{hump}} + R_{\mathrm{air}} + R_{\mathrm{surf}}\right)
\max(0.65,k_{\mathrm{surf}}).
\end{align}

Это прикладная оценка силы, которую должен преодолеть двигатель. В ней есть
динамическая часть, воздушное сопротивление и сопротивление поверхности.

\simple{Сопротивление -- это то, что мешает лодке ехать. В модели оно собрано
из нескольких частей: сопротивление корпуса, сопротивление воздуха и трение или
вязкость поверхности. Потом вся сумма усиливается коэффициентом поверхности.}
\howto{\(R_{\mathrm{dyn}}\) растет от скорости и формы движения.
\(R_{\mathrm{air}}\) -- сопротивление воздуха. \(R_{\mathrm{surf}}\) -- грубая
оценка того, насколько масса лодки давит на поверхность и насколько эта
поверхность ``тяжелая''. Финальное \(R\) складывает эти части.}
\effect{Скорость входит как \(v^2\), то есть при росте скорости сопротивление
растет очень быстро. Масса увеличивает \(R_{\mathrm{surf}}\). Плохая
поверхность увеличивает \(k_{\mathrm{surf}}\) и тоже поднимает сопротивление.}
\exampletext{Если скорость выросла в 2 раза, часть сопротивления со скоростью
может вырасти примерно в 4 раза. Поэтому быстрый режим часто заметно повышает
расход топлива.}

\subsection{Шаг 9. Требуемая мощность}

\begin{align}
P_{\mathrm{raw}} &= \frac{Rv}{1000\eta}, \\
P &= \mathrm{clamp}\left(P_{\mathrm{raw}}, P_{\min}, P_{\max}\right).
\end{align}

Мощность равна силе сопротивления, умноженной на скорость. Деление на КПД тяги
учитывает потери винта и трансмиссии.

\simple{Двигатель должен дать столько мощности, чтобы преодолеть сопротивление
на текущей скорости. Если винт и передача неидеальны, часть энергии теряется,
поэтому делим на КПД \(\eta\).}
\howto{\(R \cdot v\) -- это механическая мощность в ваттах. Деление на 1000
переводит ее в киловатты. \(\mathrm{clamp}\) ограничивает результат снизу и
сверху: модель не дает мощности стать нулевой или больше мощности двигателя.}
\effect{Больше сопротивление \(R\) -- больше мощность. Больше скорость \(v\) --
тоже больше мощность. Ниже КПД \(\eta\) -- двигатель должен работать сильнее.}
\exampletext{Если сопротивление 1800 Н, скорость 12 м/с, КПД 0.58, то
\(P_{\mathrm{raw}}\approx1800\cdot12/(1000\cdot0.58)\approx37\) кВт.}

\subsection{Шаг 10. Расход топлива}

\begin{align}
q_{\mathrm{fuel}} &=
\frac{P \cdot BSFC \cdot k_{\mathrm{mode}}}{\rho_{\mathrm{fuel}}\cdot1000}, \\
F_{\mathrm{raw}} &= q_{\mathrm{fuel}}t_{\mathrm{seg}}, \\
F_{\mathrm{fallback}} &=
d_{\mathrm{seg}}\cdot base\_l\_per\_km \cdot k_{\mathrm{surf}}
\cdot k_{\mathrm{load}}\cdot k_{\mathrm{mode}}, \\
F_{\mathrm{seg}} &=
\max\left(F_{\mathrm{raw}},\; 0.08F_{\mathrm{fallback}}\right).
\end{align}

BSFC переводит мощность двигателя в расход топлива. На выходе получается
литров в час, затем литры на участок.

\simple{Сначала считаем, сколько литров топлива уходит за час при такой
мощности. Потом умножаем на время участка и получаем топливо именно на этот
кусок маршрута.}
\howto{\(BSFC\) показывает, сколько граммов топлива нужно на 1 кВт мощности за
1 час. \(k_{\mathrm{mode}}\) делает поправку на режим: быстрый режим может быть
прожорливее, экономичный -- мягче. \(\rho_{\mathrm{fuel}}\) переводит массу
топлива в литры.}
\effect{Больше мощность \(P\), больше \(BSFC\), больше \(k_{\mathrm{mode}}\) --
расход выше. Чем дольше участок, тем больше \(F_{\mathrm{raw}}\).}
\exampletext{Если \(q_{\mathrm{fuel}}=18\) л/ч, а участок занимает 0.05 часа
(3 минуты), то \(F_{\mathrm{raw}}=18\cdot0.05=0.9\) л.}

\simple{\(F_{\mathrm{fallback}}\) -- страховочная нижняя оценка. Она нужна,
чтобы на очень коротких или странных участках расход не становился нереально
маленьким.}

\subsection{Шаг 11. Риск участка}

\begin{align}
k_{\mathrm{speed}} &=
1 + 0.22\left(\frac{V_{\mathrm{rec}}}{\max(V_{\mathrm{surface}},1)}\right)^2, \\
k_{\mathrm{narrow}} &=
\begin{cases}
1.12, & \text{узкий waterway}, \\
1.00, & \text{иначе},
\end{cases} \\
Risk_{\mathrm{seg}} &=
d_{\mathrm{seg}} r_{\mathrm{surface}} k_{\mathrm{speed}} k_{\mathrm{narrow}}.
\end{align}

Риск растет от опасной поверхности, скорости и узости русла.

\simple{Риск -- это не вероятность аварии в строгом научном смысле. Это баллы,
которые помогают алгоритму предпочитать более спокойные участки. Чем опаснее
поверхность, выше скорость и уже русло, тем больше баллов риска.}
\howto{\(k_{\mathrm{speed}}\) добавляет штраф за скорость. Скорость стоит в
квадрате, поэтому быстрый разгон поднимает риск заметнее. \(k_{\mathrm{narrow}}\)
добавляет небольшой штраф за узкие протоки и waterway-ребра.}
\effect{Большая длина участка, большой риск поверхности, высокая скорость и
узкая река увеличивают \(Risk_{\mathrm{seg}}\).}
\exampletext{Если участок длинный и идет по рисковой поверхности, безопасный
режим может выбрать путь длиннее, но с меньшим итоговым риском.}

\subsection{Шаг 12. Стоимость ребра для A*}

\begin{align}
C_{\mathrm{edge}} &=
w_d d_{\mathrm{seg}}
+ w_t\frac{t_{\mathrm{seg}}\cdot60}{10}
+ w_f\frac{F_{\mathrm{seg}}}{10}
+ w_r\frac{Risk_{\mathrm{seg}}}{5} \notag \\
&\quad
+ w_pP_{\mathrm{no\_planing}}
+ w_hH_{\mathrm{hard}}.
\end{align}

A* ищет путь с минимальной суммой \(C_{\mathrm{edge}}\).
Разные режимы маршрута отличаются весами \(w\).

\simple{Стоимость ребра -- это общий ``ценник'' участка для алгоритма. В этот
ценник входят расстояние, время, топливо, риск и штрафы. Алгоритм выбирает путь
с минимальной суммой таких ценников.}
\howto{Каждый член формулы отвечает за отдельный критерий. Вес \(w_d\) говорит,
насколько важна длина. \(w_t\) -- время. \(w_f\) -- топливо. \(w_r\) -- риск.
\(w_p\) и \(w_h\) добавляют штрафы за плохой режим движения и сложные
поверхности.}
\effect{В быстром режиме больше вес времени, поэтому алгоритм терпит больший
расход. В экономичном режиме больше вес топлива. В безопасном режиме больше вес
риска.}
\exampletext{Если два пути одинаковые по длине, но один опаснее, безопасный
режим выберет менее рисковый. Если один путь быстрее, быстрый режим может
выбрать его даже при большем расходе.}

\subsection{Шаг 13. Штрафы за потерю глиссирования и сложные поверхности}

\begin{align}
P_{\mathrm{no\_planing}} &=
\begin{cases}
0, & \mathrm{surface.planing}=\mathrm{true}, \\
d_{\mathrm{seg}}, & \mathrm{surface.planing}=\mathrm{false},
\end{cases} \\
H_{\mathrm{hard}} &=
\begin{cases}
d_{\mathrm{seg}}, & \text{сложная поверхность}, \\
0, & \text{иначе}.
\end{cases}
\end{align}

Эти штрафы помогают алгоритму избегать участков, где аэролодка идет
неэффективно или опасно.

\simple{Штраф -- это дополнительная цена, которую мы добавляем участку. Он не
обязательно означает запрет. Он просто говорит алгоритму: ``сюда лучше не ехать,
если есть нормальная альтернатива''.}
\howto{Если поверхность поддерживает глиссирование, штраф за потерю
глиссирования равен 0. Если не поддерживает, штраф равен длине участка. Для hard
поверхности штраф тоже равен длине участка.}
\effect{Чем длиннее плохой участок, тем сильнее штраф. Если конфигурация лодки
вообще запрещает сложные поверхности, такие ребра исключаются еще раньше и до
формулы стоимости не доходят.}
\exampletext{Короткий сложный участок может остаться в маршруте, если объезд
слишком большой. Длинный сложный участок алгоритм, скорее всего, постарается
обойти.}

\subsection{Шаг 14. Итоги маршрута}

\begin{align}
D &= \sum d_{\mathrm{seg}}, \\
T &= \sum t_{\mathrm{seg}}, \\
F &= \sum F_{\mathrm{seg}}, \\
Risk &= \sum Risk_{\mathrm{seg}}, \\
Fuel_{\mathrm{remain}} &= Fuel_{\mathrm{tank}} - F, \\
Fuel_{\mathrm{reserve}} &= Fuel_{\mathrm{tank}}\cdot r_{\mathrm{reserve}}.
\end{align}

После выбора маршрута движок суммирует длину, время, топливо и риск,
а затем проверяет остаток топлива относительно резерва.

\simple{На предыдущих шагах мы считали каждый маленький участок отдельно. В
конце просто складываем все участки маршрута и получаем общую длину, время,
топливо и риск.}
\howto{\(\sum\) означает ``сложить по всем сегментам''. Например,
\(F=\sum F_{\mathrm{seg}}\) значит: взять расход каждого участка и сложить в
один общий расход маршрута.}
\effect{Если маршрут стал длиннее, обычно растут \(D\), \(T\), \(F\) и риск.
Но иногда более длинный маршрут может быть безопаснее или экономичнее, если он
идет по хорошей поверхности.}
\exampletext{Если есть три участка с расходом 0.4 л, 0.7 л и 0.5 л, то общий
расход \(F=0.4+0.7+0.5=1.6\) л. Остаток топлива -- это бак минус общий расход.}

\section{Важное ограничение}

Эта модель не является точной паспортной моделью Raptor 650 и не заменяет
испытания. Она нужна для MVP и защиты кейса: показать правильную структуру
расчета и зависимость маршрута от физики лодки. Для продукта коэффициенты нужно
калибровать по реальным трекам, замерам топлива, скорости, загрузки и состояния
поверхности.

\end{document}
